\section{Estado del arte: problemas relacionados y brechas}
\label{sec:estado}

Varios estudios muestran que el crecimiento de un niño debe analizarse como un proceso y no solo mediante una medida aislada. Melo et al. \cite{r5} estudiaron la ganancia de peso en programas canguro de Bogotá; Ren et al. \cite{r6} relacionaron las trayectorias de crecimiento con factores familiares y neonatales, y Nel et al. \cite{r7} describieron patrones de recuperación y dificultades de crecimiento en 322 prematuros atendidos con el método canguro. Estos trabajos nos llevaron a considerar las visitas sucesivas como parte central del análisis. Sin embargo, no resuelven por sí mismos cómo ofrecer predicciones por visita cuando se trabaja con registros institucionales extensos, incompletos y de calidad variable.

También revisamos estudios que utilizan aprendizaje automático para anticipar dificultades del neurodesarrollo. Baker y Kandasamy \cite{r8} reúnen diferentes enfoques aplicados a prematuros, mientras que Demirci et al. \cite{r9} analizaron información de 1.109 niños muy prematuros para predecir retrasos mental y psicomotor a los 25 meses. Carlton et al. \cite{r10} estudiaron la clasificación de riesgo para seguimiento infantil, y Wang et al. \cite{r11} aplicaron modelos explicables a la hipoglucemia neonatal. Estos antecedentes confirman que los datos tempranos pueden aportar información útil, pero sus poblaciones y desenlaces son distintos de los nuestros. Por ello, no sería adecuado comparar sus cifras de rendimiento directamente con las de KMCPred. Nuestro objetivo no fue crear un algoritmo nuevo, sino integrar el análisis del riesgo con la revisión de calidad y la información que necesita el equipo clínico.

La calidad del registro también ha sido un tema central en la literatura. Kahn et al. \cite{r12} y Weiskopf y Weng \cite{r13} proponen maneras de evaluar datos de salud; Chandola et al. \cite{r14} revisan métodos para detectar anomalías. La regresión por cuantiles y la predicción conformal \cite{r15,r16,r17} ayudan a identificar valores que parecen extraños para un contexto determinado. Por su parte, la calibración y los cambios en el comportamiento de los modelos a lo largo del tiempo \cite{r18,r19,r20} son necesarios para interpretar los riesgos expresados en porcentaje. Utilizamos TreeSHAP \cite{r21} para describir qué variables influyen en las predicciones, teniendo presente que estas explicaciones muestran asociaciones, no relaciones causales \cite{r22,r23}. La guía TRIPOD+AI \cite{r24} orientó la forma de presentar la evaluación.

A partir de esa revisión, enfocamos el trabajo en tres preguntas prácticas: ¿podemos detectar registros dudosos?, ¿podemos estimar con suficiente anticipación los dos desenlaces priorizados?, y ¿podemos ofrecer estos resultados de una manera clara y reproducible? El análisis se hizo dentro de nuestro propio conjunto de datos, sin afirmar que KMCPred supera a los estudios externos \cite{r5,r6,r7,r8,r9,r10,r11}, ya que las poblaciones y los tiempos de seguimiento son diferentes.
